% 量子力学1 量子力学1-17 問題10（072-translation。当時は 量子力学1-8 問題7） 解いた日 2026-11-08
%% kotae
\toi{1}{◎}{}
$\hat{p}\ket{k}=(\hbar k)\ket{k}$ より $\hat{p}^{n}\ket{k}=(\hbar k)^{n}\ket{k}$ より
\begin{align*}
\hat{p}^{n}\hat{1}&=\int_{-\infty}^{\infty}\hat{p}^{n}\ket{k}\bra{k}\dd k\\
&=\int_{-\infty}^{\infty}(\hbar k)^{n}\ket{k}\bra{k}\dd k
\end{align*}
（脇で $\hat{T}(a)\ket{k}$ を計算）
\begin{align*}
\hat{T}(a)\ket{k}&=\sum_{n=0}^{\infty}\frac{1}{n!}\Bigl(-\frac{\ii a}{\hbar}\Bigr)^{n}\hat{p}^{n}\ket{k}\\
&=\sum_{n=0}^{\infty}\frac{1}{n!}(-\ii ak)^{n}\ket{k}\\
&=\ee^{-\ii ka}\ket{k}
\end{align*}
\begin{align*}
\hat{T}(a)\hat{1}&=\int_{-\infty}^{\infty}\hat{T}(a)\ket{k}\bra{k}\dd k\\
&=\int_{-\infty}^{\infty}\ee^{-\ii ka}\ket{k}\bra{k}\dd k
\end{align*}
\toi{2}{◎}{}
\begin{align*}
\hat{1}\ket{x}&=\int_{-\infty}^{\infty}\ket{k}\braket{k}{x}\dd k\\
&=\int_{-\infty}^{\infty}\ket{k}\frac{1}{\sqrt{2\pi}}\ee^{-\ii kx}\dd k
\end{align*}
\begin{align*}
&\hat{T}(a)\ket{x}\\
&=\sum_{n=0}^{\infty}\frac{1}{n!}\Bigl(-\frac{\ii a}{\hbar}\Bigr)^{n}\hat{p}^{n}\int_{-\infty}^{\infty}\ket{k}\frac{1}{\sqrt{2\pi}}\ee^{-\ii kx}\dd k\\
&=\int_{-\infty}^{\infty}\sum_{n=0}^{\infty}\frac{1}{n!}\Bigl(-\frac{\ii a}{\hbar}\Bigr)^{n}(\hbar k)^{n}\\
&\qquad\times\ket{k}\frac{1}{\sqrt{2\pi}}\ee^{-\ii kx}\dd k\\
&=\int_{-\infty}^{\infty}\ee^{-\ii ka}\ket{k}\frac{1}{\sqrt{2\pi}}\ee^{-\ii kx}\dd k\\
&=\int_{-\infty}^{\infty}\ket{k}\frac{1}{\sqrt{2\pi}}\ee^{-\ii k(x+a)}\dd k\\
&=\int_{-\infty}^{\infty}\ket{k}\braket{k}{x+a}\dd k\\
&=\ket{x+a}
\end{align*}
\toi{3}{◎}{}
$\hat{p}^{\dagger}=\hat{p}$
\begin{align*}
\hat{T}(a)^{\dagger}&=\Bigl(\sum_{n=0}^{\infty}\frac{1}{n!}\Bigl(-\frac{\ii a}{\hbar}\Bigr)^{n}\hat{p}^{n}\Bigr)^{\dagger}\\
&=\sum_{n=0}^{\infty}\frac{1}{n!}\Bigl(-\frac{\ii(-a)}{\hbar}\Bigr)^{n}\hat{p}^{n}\\
&=\hat{T}(-a)
\end{align*}
\begin{align*}
\hat{T}(a)^{\dagger}\hat{T}(a)&=\hat{T}(-a)\hat{T}(a)\\
&=\exp(\ii a\hat{p}/\hbar)\exp(-\ii a\hat{p}/\hbar)\\
&=\exp(\ii a\hat{p}/\hbar)\exp(\ii a\hat{p}/\hbar)^{-1}\\
&=\hat{1}
\end{align*}
\begin{align*}
\matrixel{x}{\hat{T}(a)}{\psi}&=\kakimachigai{(\hat{T}(a)^{\dagger}\ket{x})^{*}}{(\cdots)^{\dagger}}\ket{\psi}\\
&=\kakimachigai{\ket{x-a}^{*}}{\ket{x-a}^{\dagger}}\ket{\psi}\\
&=\braket{x-a}{\psi}\\
&=\psi(x-a)
\end{align*}
\toi{4}{△}{方針}
\[ \matrixel{x}{\hat{p}^{n}}{\psi}=\ayamari{\int_{-\infty}^{\infty}(\hbar k)^{n}\braket{x}{k}\braket{k}{\psi}\dd k}{(-\ii\hbar)^{n}\psi^{(n)}(x)} \]
\naze{運動量の完全性を挟んで $\hat{p}^{n}$ を固有値 $(\hbar k)^{n}$ に置きかえると，(1)〜(3) と同じ道筋になり，問われた「直接位置表示にする」確かめにならない（テイラー展開の仮定も使わない）。位置表示の $\matrixel{x}{\hat{p}^{n}}{\psi}=(-\ii\hbar)^{n}\psi^{(n)}(x)$ を使う。以降はこの道筋を引き継ぐ。}
\begin{align*}
&\matrixel{x}{\hat{T}(a)}{\psi}\\
&=\sum_{n=0}^{\infty}\frac{1}{n!}\Bigl(-\frac{\ii a}{\hbar}\Bigr)^{n}\matrixel{x}{\hat{p}^{n}}{\psi}\\
&=\sum_{n=0}^{\infty}\frac{1}{n!}\Bigl(-\frac{\ii a}{\hbar}\Bigr)^{n}\\
&\qquad\times\int_{-\infty}^{\infty}(\hbar k)^{n}\braket{x}{k}\braket{k}{\psi}\dd k\\
&=\int_{-\infty}^{\infty}\sum_{n=0}^{\infty}\frac{1}{n!}(-\ii ka)^{n}\frac{1}{\sqrt{2\pi}}\ee^{\ii kx}\tilde{\psi}(k)\dd k\\
&=\int_{-\infty}^{\infty}\frac{1}{\sqrt{2\pi}}\ee^{\ii k(x-a)}\tilde{\psi}(k)\dd k\\
&=\psi(x-a)
\end{align*}
\naze{正しくは，第 $n$ 項の係数が $(-\ii a/\hbar)^{n}(-\ii\hbar)^{n}/n!=(-a)^{n}/n!$ になり，$\sum_{n}(-a)^{n}\psi^{(n)}(x)/n!$ が $\psi(x-a)$ の $x$ のまわりのテイラー展開と項ごとに一致することで (3) を確かめる。}
\toi{5}{◎}{}
\begin{gather*}
\bra{\psi_{a}}=\kakimachigai{(\hat{T}(a)\ket{\psi})^{*}}{(\cdots)^{\dagger}}=\bra{\psi}\hat{T}(a)^{\dagger}\\
\matrixel{\psi_{a}}{\hat{x}}{\psi_{a}}=\matrixel{\psi}{\hat{T}(a)^{\dagger}\hat{x}\hat{T}(a)}{\psi}
\end{gather*}
（$\hat{x}$ と $\hat{T}(a)$ の間に $\hat{1}$ を挟む）
\begin{align*}
&=\int_{-\infty}^{\infty}\matrixel{\psi}{\hat{T}(a)^{\dagger}\hat{x}}{x}\matrixel{x}{\hat{T}(a)}{\psi}\dd x\\
&=\int_{-\infty}^{\infty}x\matrixel{x}{\hat{T}(a)}{\psi}^{*}\psi(x-a)\dd x\\
&=\int_{-\infty}^{\infty}x\braket{\psi}{x-a}\braket{x-a}{\psi}\dd x
\end{align*}
$x-a=x'$ とする。
\begin{align*}
&=\int_{-\infty}^{\infty}(x'+a)\braket{\psi}{x'}\braket{x'}{\psi}\dd x'\\
&=\int_{-\infty}^{\infty}\matrixel{\psi}{x'}{x'}\braket{x'}{\psi}\dd x'\\
&\qquad+a\bra{\psi}\Bigl(\int_{-\infty}^{\infty}\ket{x'}\bra{x'}\dd x'\Bigr)\ket{\psi}\\
&=\bra{\psi}\Bigl(\int_{-\infty}^{\infty}\hat{x}\ket{x'}\bra{x'}\dd x'\Bigr)\ket{\psi}+a\braket{\psi}{\psi}\\
&=\bra{\psi}\Bigl(\int_{-\infty}^{\infty}\hat{x}\ket{x}\bra{x}\dd x\Bigr)\ket{\psi}+a\\
&=\matrixel{\psi}{\hat{x}}{\psi}+a
\end{align*}
\[ \matrixel{\psi_{a}}{\hat{p}}{\psi_{a}}=\matrixel{\psi}{\hat{T}(a)^{\dagger}\hat{p}\hat{T}(a)}{\psi} \]
\begin{align*}
&\comm{\hat{p}}{\hat{T}(a)}\\
&=\sum_{n=0}^{\infty}\frac{1}{n!}\Bigl(-\frac{\ii a}{\hbar}\Bigr)^{n}\underbrace{\comm{\hat{p}}{\hat{p}^{n}}}_{0}=0\\
&\matrixel{\psi_{a}}{\hat{p}}{\psi_{a}}\\
&=\bra{\psi}\underbrace{\hat{T}(a)^{\dagger}\hat{T}(a)}_{\hat{1}}\hat{p}\ket{\psi}=\matrixel{\psi}{\hat{p}}{\psi}
\end{align*}
（ここまでを括弧でくくり）←これでいい？
\begin{align*}
&\matrixel{\psi_{a}}{\hat{p}}{\psi_{a}}\\
&=\int_{-\infty}^{\infty}\matrixel{\psi}{\hat{T}(a)^{\dagger}\hat{p}}{x}\matrixel{x}{\hat{T}(a)}{\psi}\dd x\\
&=\int_{-\infty}^{\infty}\matrixel{x}{\hat{p}\hat{T}(a)}{\psi}^{*}\psi(x-a)\dd x\\
&=\int_{-\infty}^{\infty}\Bigl(-\ii\hbar\frac{\dd}{\dd x}\psi(x-a)\Bigr)^{*}\psi(x-a)\dd x\\
&=\int_{-\infty}^{\infty}\ii\hbar\kakimachigai{\Bigl(\frac{\dd}{\dd x}\psi(x-a)\Bigr)}{(\cdots)^{*}}\psi(x-a)\dd x
\end{align*}
$x-a=x'$
\begin{align*}
&=\int_{-\infty}^{\infty}\ii\hbar\kakimachigai{\Bigl(\frac{\dd}{\dd x'}\psi(x')\Bigr)}{(\cdots)^{*}}\psi(x')\dd x'\\
&=\int_{-\infty}^{\infty}\matrixel{x'}{\hat{p}}{\psi}^{*}\braket{x'}{\psi}\dd x'\\
&=\int_{-\infty}^{\infty}\matrixel{\psi}{\hat{p}}{x'}\braket{x'}{\psi}\dd x'\\
&=\bra{\psi}\hat{p}\Bigl(\int_{-\infty}^{\infty}\ket{x'}\bra{x'}\dd x'\Bigr)\ket{\psi}\\
&=\matrixel{\psi}{\hat{p}}{\psi}
\end{align*}
%% kizuki
\item 位置表示では $\hat{p}$ は $-\ii\hbar\,\dd/\dd x$ なので，$\matrixel{x}{\hat{p}^{n}}{\psi}=(-\ii\hbar)^{n}\psi^{(n)}(x)$ で，$\hat{T}(a)=\exp(-\ii a\hat{p}/\hbar)$ は $\exp(-a\,\dd/\dd x)$ になる。これは問題1 (3)（047-shift）の $\exp(a\hat{D})f(x)=f(x+a)$ で $a$ を $-a$ にしたものである。運動量表示と位置表示の両方で同じ $\psi(x-a)$ が出ることが「確かめ」の中身なので，「確かめよ」と問われたら前の小問と別の道筋で計算する。
\item \mae{第3章 問題2・10-18，10-21，第5章 問題7・10-30} 複素共役とエルミート共役の書き方（減点なし）。ブラを $(\hat{T}(a)^{\dagger}\ket{x})^{*}$，$\ket{x-a}^{*}$ と書き，位置表示の計算で $(\dd\psi/\dd x)^{*}$ の $*$ を2行落とした。ケットからブラへは $\dagger$，数には $*$ と使い分け，$\abs{f}^{2}=f^{*}f$ の $*$ を落とさない。
\item \jibunmemo 交換関係の方法に「これでいい？」。この方法は正しく，解答と同じである。$\hat{T}(a)$ は $\hat{p}$ のべき級数なので $\comm{\hat{p}}{\hat{T}(a)}=0$ で，(3) の $\hat{T}(a)^{\dagger}\hat{T}(a)=\hat{1}$ と合わせれば $\hat{p}$ の期待値は変わらない。
